\textbf{Hình thức nghiên cứu:} Nghiên cứu thực nghiệm.

\textbf{Dữ liệu:} Nghiên cứu sử dụng tập dữ liệu \textbf{FoodSeg103}, gồm 7.118 ảnh được gán nhãn phân đoạn ở mức pixel cho 103 lớp nguyên liệu. Dữ liệu được chia thành 4.983 ảnh huấn luyện và 2.135 ảnh đánh giá. Các bước tiền xử lý dự kiến bao gồm chuẩn hóa kích thước ảnh, cắt ảnh ngẫu nhiên, lật ngang và biến đổi màu để tăng khả năng tổng quát hóa của mô hình.

\textbf{Thiết kế thực nghiệm:} Nghiên cứu xây dựng một pipeline thống nhất cho các mô hình CNN-based nhẹ, trong đó thay đổi có kiểm soát các thành phần chính gồm backbone, độ phân giải đầu vào và mô-đun tăng cường đặc trưng cục bộ. Cách thiết kế này nhằm tách biệt ảnh hưởng của từng yếu tố đến chất lượng phân đoạn và hiệu quả suy luận.

\textbf{Pipeline thực hiện:}
\begin{figure}[H]
    \centering
    \resizebox{\columnwidth}{!}{%
    \begin{tikzpicture}[
        node distance=0.45cm,
        >=Latex,
        every node/.style={font=\footnotesize},
        block/.style={
            rectangle,
            rounded corners=2pt,
            draw=black,
            thick,
            align=center,
            minimum height=0.85cm,
            text width=1.7cm,
            inner sep=2pt
        }
    ]
    
    \node[block] (data) {FoodSeg103};
    \node[block, right=of data] (prep) {Preprocess\\Augment};
    \node[block, right=of prep, text width=2.2cm] (model) {Lightweight\\CNN-based Model\\+ Texture-aware};
    \node[block, right=of model] (train) {Train};
    \node[block, right=of train] (eval) {Evaluate};
    \node[block, right=of eval] (analysis) {Analyze\\Trade-off};
    
    \draw[->, thick] (data) -- (prep);
    \draw[->, thick] (prep) -- (model);
    \draw[->, thick] (model) -- (train);
    \draw[->, thick] (train) -- (eval);
    \draw[->, thick] (eval) -- (analysis);
    
    \end{tikzpicture}%
    }
    \caption{Pipeline thực hiện nghiên cứu.}
    \label{fig:pipeline}
    \end{figure}

\textbf{Hướng tiếp cận mô hình:} Thay vì tập trung vào một kiến trúc cụ thể, nghiên cứu khảo sát nhóm các hướng tiếp cận CNN-based cho phân đoạn ngữ nghĩa, bao gồm cả cấu hình cơ sở và các biến thể lightweight. Trên nền các mô hình này, nghiên cứu sẽ thử nghiệm các cơ chế tăng cường đặc trưng cục bộ theo hướng texture-aware, ví dụ thông qua tinh chỉnh nhánh chi tiết, hợp nhất đặc trưng hoặc cải thiện biểu diễn biên.

\textbf{Đánh giá:} Hiệu năng mô hình được đánh giá bằng các chỉ số \textit{mean Intersection over Union} (mIoU), \textit{mean Accuracy} (mAcc), \textit{all Accuracy} (aAcc), độ trễ suy luận, kích thước mô hình và chi phí tính toán. Kết quả sẽ được phân tích theo hướng đánh đổi giữa chất lượng phân đoạn và khả năng triển khai.

\textbf{Giới hạn và tính khả thi:} Nghiên cứu giới hạn trong các mô hình CNN-based nhẹ và một benchmark chính là FoodSeg103. Việc lựa chọn phạm vi này giúp đảm bảo tính khả thi về thời gian, tài nguyên tính toán và khả năng so sánh công bằng giữa các mô hình.
